\chapter{Introduction générale}

Au cours des dernières années, la transformation numérique est devenue un levier stratégique pour les organisations publiques et privées. Cette évolution s'accompagne d'une croissance exponentielle du volume de données produites quotidiennement par les différents systèmes d'information. Dans le domaine de la santé, cette quantité importante de données constitue une ressource précieuse pour le suivi des activités hospitalières, l'évaluation des performances, la recherche scientifique ainsi que l'aide à la prise de décision.

Parallèlement, les progrès réalisés dans le domaine de l'intelligence artificielle, notamment avec l'émergence des grands modèles de langage (\textit{Large Language Models -- LLM}), ont profondément modifié les modes d'interaction entre l'utilisateur et les systèmes informatiques. Ces modèles permettent aujourd'hui de comprendre des requêtes formulées en langage naturel, de raisonner sur des informations complexes et de générer des réponses cohérentes tout en facilitant l'accès à des connaissances auparavant réservées à des utilisateurs spécialisés.

Malgré ces avancées, l'analyse de données demeure une tâche relativement complexe. Les outils traditionnels nécessitent généralement la maîtrise de langages de programmation tels que Python ou R, de bibliothèques spécialisées comme \texttt{pandas}, ou encore de plateformes décisionnelles dédiées. Cette dépendance aux compétences techniques constitue un frein important pour de nombreux utilisateurs métiers souhaitant exploiter leurs données de manière autonome.

% \begin{infobox}{Contexte institutionnel}
Dans ce contexte, le \textbf{Groupement Sanitaire Territorial Tanger--Tétouan--Al Hoceïma (GST-TTA)} poursuit sa stratégie de transformation numérique afin d'améliorer l'exploitation des données issues de ses établissements de santé. Cette démarche vise à développer des solutions innovantes permettant de faciliter l'accès à l'information, d'améliorer les capacités d'analyse et de soutenir la prise de décision.

C'est dans cette perspective que s'inscrit ce projet de fin d'année, dont l'objectif est de concevoir une \textbf{plateforme intelligente d'analyse de données conversationnelle}. La solution développée permet aux utilisateurs d'interagir avec leurs jeux de données en langage naturel, de réaliser automatiquement des analyses statistiques, de générer des visualisations pertinentes et d'obtenir des rapports synthétiques grâce à l'intégration d'agents d'intelligence artificielle.
% \end{infobox}

\vspace{0.5cm}

\section{Organisation du rapport}

Le présent mémoire est organisé en six chapitres complémentaires, suivis d'une conclusion générale.

\begin{table}[H]
\centering
\caption{Structure du mémoire}
\label{tab:structure}
\begin{tabularx}{\textwidth}{>{\bfseries}p{4.5cm}Y}
\toprule
\textbf{Partie} & \textbf{Contenu} \\
\midrule
Chapitre 1 --- Introduction générale & Situe le contexte de la transformation numérique du secteur de la santé, présente la problématique et l'organisation du mémoire. \\
\midrule
Chapitre 2 --- Contexte général & Présente l'organisme d'accueil (GST-TTA), l'analyse de l'existant, la problématique, les objectifs et la méthodologie adoptée. \\
\midrule
Chapitre 3 --- Analyse des besoins & Décrit les besoins fonctionnels et non fonctionnels, les acteurs, l'état de l'art, l'étude comparative et le positionnement de la solution. \\
\midrule
Chapitre 4 --- Conception et architecture & Détaille la conception générale, la modélisation UML, l'architecture logicielle, l'architecture agentique, le moteur d'analyse et la gestion des données. \\
\midrule
Chapitre 5 --- Réalisation & Présente la mise en œuvre effective : backend, orchestration et evidence, moteur d'analyse, interface et internationalisation, déploiement et infrastructure de modèles locaux. \\
\midrule
Chapitre 6 --- Tests et validation & Expose la stratégie de test, les scénarios réalisés, les résultats obtenus, les limites et les perspectives d'amélioration. \\
\midrule
Conclusion générale & Synthétise les travaux réalisés, les apports du projet et les perspectives d'évolution. \\
\bottomrule
\end{tabularx}
\end{table}

\begin{resultbox}{Synthèse}
Cette organisation permet de présenter progressivement les différentes étapes ayant conduit à la conception et au développement de la plateforme tout en mettant en évidence les contributions techniques et scientifiques apportées par ce projet.
\end{resultbox}